\documentclass[12pt,a4paper]{exam}

% --------------------
% Packages
% --------------------
\usepackage[utf8]{inputenc}
\usepackage[T1]{fontenc}
\usepackage[french]{babel}
\usepackage{amsmath,amssymb}
\usepackage{multicol}
\usepackage{geometry}

\geometry{margin=2cm}

% --------------------
% En-tête / pied
% --------------------
\header{}{}{2nde générale}
\footer{}{\thepage}{\it{G.M.}}

% --------------------
\renewcommand{\questionlabel}{\alph{question})}
\begin{document}

\begin{center}
    \Large\textbf{Exercices demi groupe\\Équations et inéquations}
\end{center}

\vspace{5pt}

% ======================================================
\noindent
\textbf{Exercice 1 : Équations}

\vspace{5pt}
Résoudre les équations suivantes dans $\mathbb{R}$ :
\vspace{-5pt}
\begin{questions}
\begin{multicols}{2}
    \question $3x - 5 = 7x$
    \question $-\dfrac{2}{3}x + 4 = 10$
\end{multicols}
\vspace{5pt}
\begin{multicols}{2}
    \question $(x-2)(x+5)=0$
    \question $x(3x-1)=0$
\end{multicols}
\vspace{5pt}
\begin{multicols}{2}
    \question $x^2 -3 = 0$
    \question $x^2 + 7= 0$
\end{multicols}
\vspace{5pt}
\begin{multicols}{2}
    \question $\displaystyle \frac{3x-4}{x+1}=0$
    \question $\displaystyle \frac{2x}{-x-3}=0$
\end{multicols}
\end{questions}

\vspace{10pt}

% ======================================================
\noindent
\textbf{Exercice 2 : Inéquations}

\vspace{5pt}
Résoudre les équations suivantes dans $\mathbb{R}$ :
\vspace{-5pt}
\begin{questions}
\begin{multicols}{2}
    \question $5x - 7 \ge 3$
    \question $-3x + 2 < 6$
\end{multicols}
\vspace{5pt}
\begin{multicols}{2}
    \question $(x-1)(x+3)\ge 0$
    \question $x(2-x)<0$
\end{multicols}
\vspace{5pt}
\begin{multicols}{2}
    \question $\displaystyle \frac{x-1}{x+2} \le 0$
    \question $\displaystyle \frac{2x+1}{x-4} > 0$
\end{multicols}
\end{questions}

\vspace{10pt}

% ======================================================
\noindent
\textbf{Exercice 3 : Étude de signe d’une fonction}

\vspace{5pt}
Étudier le signe des fonctions suivantes sur leur ensemble de définition :
\vspace{-5pt}
\begin{questions}
\begin{multicols}{2}
    \question $f(x)=(x-2)(x+4)$
    \question $g(x)=-x(3x-5)$
\end{multicols}
\vspace{5pt}
\begin{multicols}{2}
    \question $\displaystyle h(x)=\frac{x-3}{x+1}$
    \question $\displaystyle k(x)=\frac{2x}{x-5}$
\end{multicols}
\end{questions}

\vspace{10pt}

% ======================================================
\noindent
\textbf{Exercice 4 : Approfondissement}

\vspace{5pt}
Résoudre les équations suivantes dans $\mathbb{R}$ :
\vspace{-5pt}
\begin{questions}
\begin{multicols}{2}
    \question $(6x+4)(8x-5) = (9x+6)(4x-1)$
    \question $-3x^2+18x-27=0$
\end{multicols}
\vspace{5pt}
\begin{multicols}{2}
    \question $\dfrac{4x-1}{6x+9} = \dfrac{2x+5}{3x-11}$
    \question $\dfrac{9x^2-42x+49}{9x^2-49} = 0$
\end{multicols}
\end{questions}

\end{document}
